\begin{definition} \label{definition:absolute_continuity_of_a_real_valued_function_on_an_interval}
    \TODO{AI generated, verify}
    Let $I \subseteq \mathbb{R}$ be an interval and let $f : I \to \mathbb{R}$ be a \CrefAndHyperrefIfExist{definition:function_of_sets}{function}.

    The function $f$ is called \hldef{absolutely continuous on $I$} if for every $\varepsilon > 0$ there exists $\delta > 0$ such that for any finite collection of pairwise disjoint subintervals $(a_1, b_1), \dots, (a_n, b_n)$ of $I$ satisfying
    $$\sum_{k=1}^n (b_k - a_k) < \delta,$$
    we have
    $$\sum_{k=1}^n |f(b_k) - f(a_k)| < \varepsilon.$$

    Absolute continuity implies \CrefAndHyperrefIfExist{definition:uniformly_continuous_function_between_extended_metric_spaces_valued_in_ordered_commutative_monoid_with_adjoined_infinity}{uniform continuity} and \CrefAndHyperrefIfExist{definition:function_of_bounded_variation_on_an_interval}{bounded variation}. On a closed bounded interval, a function is absolutely continuous if and only if it is \CrefAndHyperrefIfExist{definition:derivative_of_a_real_valued_function_on_a_real_interval}{differentiable} almost everywhere, its derivative is \CrefAndHyperrefIfExist{definition:lebesgue_integral_of_measurable_functions_into_the}{Lebesgue integrable}, and the \CrefAndHyperrefIfExist{theorem:fundamental_theorem_of_calculus_part_one_antiderivative_from_integral}{fundamental theorem of calculus} holds:
    $$f(x) - f(a) = \int_a^x f'(t)\,dt$$
    for all $x \in [a, b]$.
\end{definition}